% ---
% id: meat-core/frank-wolfe/rem:fo-arithmetic
% kind: rem
% title: "Exact real arithmetic in the first-order bound"
% summary: >-
%   Both linear minima are bounded exactly over the exact sets, but for the
%   gradient as computed: the objective, the gradient and the maps and frame
%   behind them are binary64, and a reference's centres move the measured
%   family. So a computed h_lb may exceed h_low, and the zero clamp repairs
%   neither.
% prominence: standard
% status: accepted
% uses: [thm:firstorder, cor:fo-keylength, cor:fo-reference, lem:feasible-side-dual, meat-core/prop:optol-nobound, meat-core/rem:optol-direction]
% mentions: [req:certified-evaluation, lem:certified-bracket-enclosed, rem:certified-scope, foundations/ass:snapped-measurement, req:iterate-checks, req:bound-estimate-difference, req:bracket-on-reference, req:verified-feasible-side, rem:acc-side, rem:fo-check-exposures, meat-core/thm:choi, meat-core/req:companion-admission, foundations/req:exact-island, foundations/def:exact-island, foundations/def:tiers, foundations/lem:recentred-enclosure, reductions/lem:row-relaxation, reductions/req:reference-companion]
% assumes: []
% refs: []
% provenance: newly-derived
% lean:
% ---

\begin{remark}\label{rem:fo-arithmetic}
Theorem~\ref{thm:firstorder}\,(ii), respectively~(iiib) --- the part a
run is in, its iterates being established only in the relaxed set
(Requirement~\ref{req:iterate-checks}) --- with the lower bounds of its
part~(iv) in place of both pairings, which is how a run forms every bracket
(Remark~\ref{rem:acc-side}, Requirement~\ref{req:verified-feasible-side}),
concludes $\hlb\le\hlow$ in exact real arithmetic, for the exact objective
and the exact gradient. No assumption about a solver enters it. What the
theorem still asks, and a run cannot verify, is exact real arithmetic in
the terms a run computes in binary64. What that means for a number a run
prints is narrower than the theorem, and a reader who takes \eqref{eq:hlb}
for a runtime invariant is wrong in the unsafe direction.

\emph{What is exact, and of what.} The two lower bounds that stand in for
the minima are certified exactly
(\xref{foundations/def:exact-island}\,(ii), (iii)), and both are bounds over
the exact sets, but of the linear form at the gradient as the run computed
it. The acceptance-side one is the exact minimum over $\acSet$ of that
linear form. The feasible-side one is at most the minimum over the exact
$\Feas^{J}$ of that linear form: the run builds its constraint rows with a
binary64 frame $\widetilde W_{\mathsf t}$ in place of the frame
$W_{\mathsf t}$ of Theorem~\ref{thm:choi}, which carries a square root of
the declared marginal, and Lemma~\ref{lem:feasible-side-dual}\,(iv) charges
the difference, a deduction linear in the row multipliers and in a bound on
$\opn{\widetilde W_{\mathsf t}-W_{\mathsf t}}$; where no such bound is
formed the row multipliers are zero and the rows are not read. So no status
of the backend, no dual it returns and no rounding of the frame can raise
the feasible-side bound above that minimum. Without the deduction it could
be raised: Lemma~\ref{lem:feasible-side-dual}\,(i) applied to rows built
from $\widetilde W_{\mathsf t}$ bounds the minimum over the set those rows
cut out, which need not contain $\Feas^{J}$ and may be empty, and over an
empty set every number is a lower bound.

\emph{The conclusion is not a floating-point invariant.} The objective
$O^{J}_{\alpha}(x_t)$, the gradient $(\lmoq_t,\lmoJ_t)$ and the subtracted
pairings $\langle\lmoq_t,\mathbf q_t\rangle$ and $\langle\lmoJ_t,J_t\rangle$
are computed in binary64 (\xref{foundations/req:exact-island}), and so are
the data the engine derives from the declaration for them: the
\tref{statistics map}, the \tref[branch map]{branch maps} and their supports, and the frame
$\widetilde W_{\mathsf t}$
through which they read Alice's side. The gap between the gradient as
computed and the exact gradient is part of that unbudgeted evaluation
arithmetic, and no claim of this module bounds it. A computed $\hlb$ may
therefore exceed the computed estimate, and may exceed the true $\hlow$, by
an amount no claim of this module bounds. Two consequences follow, and
neither is a defect of the theorem.
\begin{enumerate}[label=(\alph*),leftmargin=2.2em]
\item \emph{$\max\{0,\hlb\}$ does not repair it.} The clamp is sound and
  it is what Corollary~\ref{cor:fo-keylength} consumes, but it corrects
  only the \emph{low} side: clamping at zero cannot bring down a bracket
  that came out too high. The clamp is no general safety net, and a run is
  not to be written, or read, as if it were.
\item \emph{A comparison of $\hlb$ against the reported estimate is a
  diagnostic, never an assertion.} The ordering of $\hlb$ below the
  objective at the returned point is a consequence of exact arithmetic and
  of the returned point's exact membership in $\acSet\times\Feas^{J}$
  (Theorem~\ref{thm:firstorder}\,(i) and~(ii)), neither of which this module
  establishes for a run: its returned point is established at most in the
  relaxed set, where the objective is bounded below only by
  $\hlow[\Feas^{J}(\varepsilon)]\le\hlow$
  (Theorem~\ref{thm:firstorder}\,(iiia)), and near convergence, where
  $\hlb$ is close to $\hlow$, that alone can put the estimate below $\hlb$.
  The difference $\hlb$ minus the estimate is recorded with its sign
  (Requirement~\ref{req:bound-estimate-difference}), and a positive
  difference, a bound above the estimate, is evidence about the arithmetic
  or about the relaxation, and not by itself evidence that the first-order
  tier is unsound.
\end{enumerate}

\emph{Evaluation on a reference companion.} Where the bound is formed on a
compiled reference companion $\Decl^{\mathrm{ref}}$
(Corollary~\ref{cor:fo-reference},
Requirement~\ref{req:bracket-on-reference}), the theorem is applied to the
companion, a declaration of the reference model, but what is evaluated is
not that declaration. Its measured family is known through an entrywise
enclosure, recentred so that every entry whose enclosure contains zero has
centre zero (\xref{foundations/lem:recentred-enclosure}), and it is
evaluated at the binary64 numbers nearest to the centres. The evaluated
family differs from the companion's, entry by entry, by at most the radius
and one rounding, and it need not be a POVM: its completeness defect and its
least eigenvalue are measured and recorded, never decided
(Requirement~\ref{req:companion-admission}\,(iii)). The objective, the
gradient, the \tref[branch map]{branch maps} and their supports, the key-inert pairs and the
rank comparison behind (H-SUPP) are all formed from those numbers, and so is
the linear form whose minimum the feasible-side bound bounds. Recentring is
applied to the handed
blocks and again to the companion's datum after the construction has
combined their entries (\xref{reductions/req:reference-companion}\,(i)), so
that an entry of the datum whose exact value is a sum that cancels to zero
is not left with a centre that is not zero. After it, a centre is not zero
only where its enclosure excludes zero; the compiler's weighting multiplies
a centre by a weight and its radius by the weight's modulus, and a product of
factor families multiplies single entries, so the compiled family keeps that
property, and rounding a centre to binary64 can turn it to zero but never
the reverse. The evaluated family therefore has a non-zero entry only where
the companion's family has one; it may lack one where the companion's entry
is no larger than the recentred radius there, and a support,
or a key-active pair, that the companion has through such entries alone is
not seen by the evaluation. What the theorem asks of the point itself,
(H-SUPP), holds of the companion whenever every block of $J_t$ is positive
definite (Corollary~\ref{cor:fo-reference}\,(ii)). This exposure is
larger in kind than the binary64 evaluation of the objective and the
gradient on any declaration. There the inputs are exact and only the
formulas are evaluated in binary64. Here the measured family itself is
displaced, by up to the radius and one rounding per entry, with a relative
error up to order one on an entry not much larger than its radius. The
displacement from the companion to its evaluation is a perturbation of the
measured family, and Proposition~\ref{prop:optol-nobound} applies to it:
closeness of two measured families bounds the change of the value by
nothing, and the change may point either way
(Remark~\ref{rem:optol-direction}); the evaluated family need not even be a
POVM. Its size is set by the radius of the enclosure, and not by binary64
rounding alone, save on an exact reference, whose radius is zero. No claim
of this module bounds its effect, and the largest radius, which the record
carries, is the only handle on it. Evaluating on the companion removes the
premise that the device measures the snapped family
(\xref{foundations/ass:snapped-measurement}); it does not remove this
exposure, which narrows only with the enclosure. The companion's entries are
not part of this exposure: they are emitted in binary64 as the oracle reads
them, their rounding absorbed into their constants
(\xref{reductions/lem:row-relaxation}), so they are the companion's own
exact rows, and the feasible-side bound is over the companion's exact
$\Feas^{J}$, the frame's rounding charged as on any declaration. The
first-order number is not evaluated over the enclosure; the certified number
of Requirement~\ref{req:certified-evaluation} is, where every element the
gradient reads has a rank fixed throughout the enclosure
(Lemma~\ref{lem:certified-bracket-enclosed}), and it then holds of the
companion itself.

\emph{Where the exposure is.} Of the ingredients of \eqref{eq:hlb}, the
objective and the gradient are evaluated by formulas this module states
in closed form, at points where (H-POS) keeps every denominator bounded
away from zero; the subtracted pairings are inner products; and the two
lower bounds are certifications of the \tref{exact island} over the exact
sets, for the gradient as computed. No ingredient is a solver's return: the
oracle's point and its status are diagnostics
(Requirement~\ref{req:verified-feasible-side}), and a solve that stopped
short of its optimum, or failed, loosens a bracket and cannot raise it
above the value it has at the exact minima of the computed linear forms.
What a run leaves unverified is named: the binary64 evaluation of the
objective, of the gradient and of the maps and frame they are built from,
of the subtracted pairings, and of the checks whose one-sided exposures
Remark~\ref{rem:fo-check-exposures} names; and on a reference companion
the displacement to its centres. It is \emph{unbudgeted}: no claim of this
module bounds it, and no assumption of this module covers it, the engine
computing outside the exact island without an error budget, as
\xref{foundations/req:exact-island} allows. That is what the
\tref{first-order} tier costs: it is weaker than a verified numerical
certificate of $\hlow$, a first-order number being a lower bound
established under the hypotheses that the claim stating it names
(\xref{foundations/def:tiers}), evaluated in arithmetic that no claim of this
module budgets. The \tref{certified} tier of the same definition is the bound
formed again with that arithmetic exact or enclosed, and what it still rests
on is Remark~\ref{rem:certified-scope}.
\end{remark}
